\chapter{Qualche esempio}

\section{Come scrivere il testo.}

Vai a capo dopo ogni punto.
E qui scrivi la frase successiva.

Vai a capo due volte se vuoi iniziare un nuovo paragrafo.

Try to \emph{empathise words} or make them \textbf{bold}.

\section{Example lists}

Non numerated lists:
\begin{itemize}
  \item list entries start with the \texttt{item} command;
  \item individual entries are indicated with a black dot, a so-called bullet;
  \item the text in the entries may be of any length.
\end{itemize}

Numerated list:
\begin{enumerate}
  \item items are numbered automatically;
  \item the numbers start at 1 with each use of the \texttt{enumerate} environment;
  \item another entry in the list.
\end{enumerate}

Una lista nel testo:
\begin{itemize*}
    \item una lista in linea;
    \item un altro elemento;
    \item un altro ancora.
\end{itemize*}

Una lista con lettere maiuscole:
\begin{enumerate*}[label=\Alph*)]
    \item una lista in linea;
    \item un altro elemento;
    \item un altro ancora.
\end{enumerate*}

Una lista con lettere minuscole:
\begin{enumerate*}[label=\alph*.]
    \item una lista in linea;
    \item un altro elemento;
    \item un altro ancora.
\end{enumerate*}

\section{Example figure and table}

Una semplice tabella inserita nel testo.
Il testo, a meno che non si tratti di intestazioni, dovrebbe essere sempre allineato a sinistra.
Evita di utilizzare le linee verticali ed utilizza solo quelle orizzontali.

% @{} c c @{} indica come sono strutturate le colonne
% @{}: rimuove gli spazi sul bordo della tabella
% l (elle): indica la presenza di una colonna nel testo il cui testo è centrato
% puoi allineare il testo al centro con la colonna 'c', o a destra con la colonna 'r'
\begin{tabular}{@{} l l @{}}
\toprule
    intestazione & intestazione \\
\midrule
    corpo & corpo \\
\bottomrule
\end{tabular}

% Posizione nello specificatore float facoltativo
% h: here (qui)
% t: top (all'inizio della pagina)
% b: bottom (alla fine della pagina)
% p: page (una pagina dedicata a questa tabella)
\begin{table}[h]
    \centering
    % inserisci la didascalia *sopra* la tabella e non lasciare uno spazio fra 'caption' e 'label'
    \caption{Caption}\label{tab:my_label}
    \begin{tabular}{@{} *{2}{l} @{}}
    \toprule
        intestazione & intestazione \\
    \midrule
        corpo & corpo \\
    \bottomrule
    \end{tabular}
\end{table}

\section{Example figure and table with the tabularray package}

Attenzione il pacchetto \package{tabularray} seppur molto consuma prende molto tempo di compilazione, è sconsigliato utilizzarlo su Overleaf.
Per questo motivo è presente il codice \verb|\IfPackageLoadedTF{tabularray}| il quale controlla se il paccheto è stato caricato e in caso affermativo esegue il codice, altrimento stampa semplicemente un avvertimento.

\IfPackageLoadedTF{tabularray}{
La tabella~\ref{tab:tabella-con-tabularray}

\begin{table}[h!tbp]% here top bottom page
    \centering% allineata al centro
    \caption{Caption}% <- importante, non rimuovere questo '%'
    \label{tab:tabella-con-tabularray}% tab convenzione per 'table'
    \begin{tblr}{
    % opzioni della tabella
        % colspec={ *{3}{Q[c,m]} },% 3 colonne, il cui contenuto è centrato orizzontalmente e verticalmente
        row{1}={font=\bfseries},
        hline{1,Z} = {2pt},% prima ed ultima riga
        hline{2} = {1pt},
        hline{3-Y} = {.1pt, lightgray},% Y penultima riga
    }
    % contenuto della tabella
    Alpha & Beta & Gamma & Delta \\
    Epsilon & Zeta & Eta & Theta \\
    Iota & Kappa & Lambda & Mu \\
    Nu & Xi & Omicron & Pi \\
    Rho & Sigma & Tau & Upsilon \\
    Phi & Chi & Psi & Omega \\
    \end{tblr}
\end{table}

\begin{tblr}{
    colspec={ Q[l] Q[c] Q[r] },% allineamento colonne: (l)eft, (c)enter, (r)ight
    rowspec={ Q[t] Q[m] Q[b] },% allineamento right: (t)op, (m)iddle, (b)ottom
    hlines, vlines,% for debug purposes only
}
    {Alpha \\ Alpha} & Beta & Gamma \\% puoi spezzare il contenuto di una cella usando le '{}'
     Delta & Epsilon & {Zeta \\ Zeta} \\
     Eta & {Theta \\ Theta} & Iota \\
\end{tblr}

Esempio di tabella che prende tutta la pagina:

\begin{tblr}{
    colspec={@{} XXX @{}},
    hline{1,Z},% prima ed ultima riga
    vlines,% for debug purposes only
}
    Alpha & Beta & Delta\\
    Gamma & Delta & Epsilon\\
    Zeta & Eta & Theta\\
\end{tblr}

\begin{tblr}{
    hlines = {1pt,white},
    row{odd} = {blue7},
    row{even} = {azure7},
    column{1} = {purple7,c},
}
    Alpha & Beta & Gamma & Delta \\
    Epsilon & Zeta & Eta & Theta \\
    Iota & Kappa & Lambda & Mu \\
    Nu & Xi & Omicron & Pi \\
    Rho & Sigma & Tau & Upsilon \\
    Phi & Chi & Psi & Omega \\
\end{tblr}

\begin{tblr}{
    row{odd} = {bg=lightgray},
    % l'ordine dei comandi è importante
    row{1} = {bg=darkgray, fg=white, font=\sffamily},
}
    Alpha & Beta & Gamma \\
    Delta & Epsilon & Zeta \\
    Eta & Theta & Iota \\
    Kappa & Lambda & Mu \\
    Nu Xi Omicron & Pi Rho Sigma & Tau Upsilon Phi \\
\end{tblr}

\(\begin{tblr}{
    column{1} = {mode=text},
    column{3} = {mode=dmath},% display math
}
    Alpha   & \frac{1}{2} & \frac{1}{2} \\
    Epsilon & \frac{3}{4} & \frac{3}{4} \\
    Iota    & \frac{5}{6} & \frac{5}{6} \\
\end{tblr}\)
}{Il pacchetto \texttt{tabularray} non è stato caricato.}

\section{Wrap text of tables in text}

\IfPackageLoadedTF{wrapfig2}{
\begin{wrapfigure}{r}{50mm}
\centering\unitlength=1mm
\begin{picture}(40,30)
\polyline(0,0)(0,30)(40,0)(0,0)(40,30)(0,30)
\Line(40,0)(40,30)
\end{picture}
\caption{A rectangle with its diagonals}\label{fig:figure}
\end{wrapfigure}
{\itshape \kant[1]}
}{Il pacchetto \texttt{wrapfig2} non è stato caricato.}

For more example take a look at the package documentation of \package{wrapfig2}.

\section{Inserimento del codice}

Attenzione il pacchetto \package{minted} seppur molto consuma prende molto tempo di compilazione, è sconsigliato utilizzarlo su Overleaf.
Per questo motivo è presente il codice \verb|\IfPackageLoadedTF{minted}| il quale controlla se il paccheto è stato caricato e in caso affermativo esegue il codice, altrimento stampa semplicemente un avvertimento.

% \IfPackageLoadedTF{minted}{
% Inserimento di uno spezzone breve di codice direttamente nel testo:
% \begin{minted}{java}
% public class Hello {
%    public static void main(String[] args) {
%       System.out.println("Hello, World!");
%    }
% }
% \end{minted}
% }{Il pacchetto \texttt{minted} non è stato caricato.}